% Source: https://artofproblemsolving.com/community/c3h1647999_09_really_equals_1
% Stable identity: aops-topic-1647999-q01
\textbf{(a)} Let $s_n=0.\underbrace{9\cdots9}_{n}=\dfrac{9}{10}+\dfrac{9}{10^2}+\cdots+\dfrac{9}{10^n}$. Then $1-s_n=\dfrac{1}{10^n}$, which tends to $0$, so $s_n\to1$ and $0.999\ldots=1$.

\textbf{(b)} The decimal is $u+0.\underbrace{0\cdots0}_{k}999\ldots$, and the second part is $10^{-k}\cdot0.999\ldots=10^{-k}$ by (a). So it equals $u+10^{-k}$, the decimal $u$ with its last digit increased from $d$ to $d+1$ (no carrying is needed, since $d\le8$). For instance $5.3999\ldots=5.3+0.0999\ldots=5.3+0.1=5.4$.

\textbf{(c)} ``$0.000\ldots1$'' is not a decimal: a decimal has one digit in each of the positions $1,2,3,\dots$ after the point, and if all of them are $0$ there is no position left for the $1$. So the first equation has no meaning, and ``we never get to the $1$'' is not a proof of anything. What is true, and is the correct form of the idea, is the statement about finite stages: $1-s_n=10^{-n}$, and these differences become smaller than any positive number. The only real number $x\ge0$ with $x\le10^{-n}$ for every $n$ is $0$ (if $x>0$, choose $n$ with $10^n>\dfrac{1}{x}$). Since $0\le1-0.999\ldots\le1-s_n=10^{-n}$ for every $n$, we get $1-0.999\ldots=0$, as in (a).
